% Section of the radiation safety manual in Spanish
\section*{Aislamiento y manejo de fuentes radiactivas dañadas}

% Itemized instructions for isolating radiation sources
\begin{itemize}
    \item[7.7.8] Con cuidado vaya aislando las diferentes fuentes, alejándolas del lugar hasta encontrar la o las causantes de los niveles de radiación.
    \item[7.7.9] Si la procedencia es de algún contenedor dañado, prepare un baño grande con agua (aproximadamente de 35-45 cm de profundidad).
    \item[7.7.10] Vea la manera de introducir el contenedor en el baño. Esto disminuirá el nivel de radiación.
    \item[7.7.11] Vea si puede retirar los tapones del equipo. Utilice las herramientas que sean necesarias.
    \item[7.7.12] Una vez hecho esto conecte la fuente a un mineral y vea si puede retirarla del contenedor dañado.
    \item[7.7.13] Si esto es posible trasladelo a un contenedor seguro.
    \item[7.7.14] Si la fuente no puede ser retirada del equipo, abandone la tarea y prepare un blindaje adecuado para todo el contenedor. Trasládelo como desecho al ININ, o si es posible, al proveedor.
\end{itemize}
